% Created 2026-09-18 Fri 14:47
% Intended LaTeX compiler: pdflatex
\documentclass[11pt]{article}
\usepackage[utf8]{inputenc}
\usepackage[T1]{fontenc}
\usepackage{graphicx}
\usepackage{longtable}
\usepackage{wrapfig}
\usepackage{rotating}
\usepackage[normalem]{ulem}
\usepackage{amsmath}
\usepackage{amssymb}
\usepackage{capt-of}
\usepackage{hyperref}
\author{fcalado}
\date{\today}
\title{}
\hypersetup{
 pdfauthor={fcalado},
 pdftitle={},
 pdfkeywords={},
 pdfsubject={},
 pdfcreator={Emacs 29.3 (Org mode 9.6.15)}, 
 pdflang={English}}
\begin{document}

\tableofcontents

\section{Research Statement}
\label{sec:org0134cf7}
My work uses digital tools and methodologies to study language related
to gender, sex, and sexulaity. While my formal training is in English
Literature, with subfocuses in Queer Studies and the Digital
Humanities (DH), I work at the intersection of the Humanities and
Computer Science disciplines, on topics that span Trans studies, AI
ethics, and digital scholarly editing.

Since beginning my position at the Pratt Institute School of
Information in Fall 2024, I pursue a research agenda based on three
main goals:

(1) to interrogate the contradictions of present-day theorizations of
gender, sex, and sexuality; 

(2) to read computational mechanisms as heuristics for cultural norms;
and

(3) to communicate across the Humanities and Computer Science
disciplines.

As of Fall 2026, I have written five papers, two that have been
published in peer-reviewed venues, and three that are currently being
considered for publication at peer-reviewed venues. In addition, I
have also secured an academic book contract, with a manuscript due in
spring 2027.

In what follows, I explain the details of this work and how they fit
into the goals described above. 

\subsubsection{To Interrogate The Contradictions Of Present-Day Theorizations Of In Gender, Sex, And Sexuality}
\label{sec:orgb92caae}
Much of my more recent research takes contemporary discourses on
gender, sex, and sexuality as its primary subject. The aim of my
research is to disentangle some of the contradictions that constitute
social and cultural definitions of these terms. Over the past decade,
the discourse on gender in particular has been highly polarized, with
various gender movements, including reactionary movements of what is
called "gender-critical feminism" and "TERFism," which are driven by
conservative arguments opposing "gender ideology" by grounding
defintions of gender in "biological" sex.

I explore these contradictions in two recent papers, both which are
currently being considered at peer-reviewed venues.
\begin{itemize}
\item The first paper, "What LLMs Can't See About Gender," under review at
\emph{FAccT 2027}, argues that some gender bias evaluation methods
effectively limit how LLMs may detect and mitigate gender bias.
These methods, which define gender as primarily an internal
phenomenon, miss the ways that geneder is expressed as external or
behavioral in language. Drawing from recent theorizing in Trans
Studies theorizing on the subject, the paper argues that defining
gender as internal only reinscribes dominant, normative gender and
furthers marginalizing visibly non-normative genders.
\item the paper explains that when gender is defined as internal, genders
which are apparently normative will be read as "default", while
non-normative genders become hyper-visible and therefore vulnerable
to stigmatization.
\item The second paper, "'My Body Feels Like it’s Coming Off': on Cis
Embodiment in \emph{Love Is Blind}," under review at \emph{TSQ: Trans Studies
Quarterly}, draws from cisgender experiences of disembodiment to
theorize common ground between trans and cis experiences of bodily
dissonance. It analyzes a dating reality TV show, \emph{Love Is Blind},
which places its cisgender participants into an environment where
they cannot see their partners, which evokes aspects of bodily
dissociation as theorized in Trans Studies. The paper argues that
these experiences create a common ground between cis and trans
relationships to their own bodies, even despite their dramatically
different life experiences and identities.
\end{itemize}

\subsubsection{Reading Computational Mechanism As Heuristics For Cultural Norms}
\label{sec:org0116076}

This research theme is the main argument of my forthcoming book,
\emph{Queer Code: Re-working Language Technologies}, which is in contract
with Bloomsbury Academic Press.

The book explores how language technologies, like web scraping, text
analysis, and generative AI, reduce language complexity to make it
computable. In these technologies, words are removed from their
original context, cleaned and standardized, so that they can be stored
and analyzed at scale. These increasingly automated processes encode
dominant logics and extractive incentives onto language data, creating
well-documented issues with data sovereignty, privacy, and most
recently, bias, as current AI tools demonstrate.

Each of the book's five chapters considers how these processes of
reduction operate through different language technologies, identifying
a mechanism in each technology that reduces or flattens language
meaning. Then, drawing from Queer Studies and related fields, the book
identifies conceptual tools that re-work these constraints into new
analytical methodologies to open up novel readings of language data.

The table below outlines the technology, constraint, and conceptual
tool across each of the book's five chapters.

\begin{center}
\begin{tabular}{llll}
Chapter & Technology & Mechanism & Concept\\[0pt]
\hline
1. Text Capture & Web Scraping & Protocol & Reparative Reading\footnotemark\\[0pt]
2. Text Analysis & Text Analysis & Iteration & Gender Performativity\footnotemark\\[0pt]
3. Text Preservation & Markup Languages & Containment & Dis-Identification\footnotemark\\[0pt]
4. Text Display & Screen Display & Abstraction & The Flesh\footnotemark\\[0pt]
5. Text Generation & Machine Learning & Aggregation & Body Dysphoria\footnotemark\\[0pt]
\end{tabular}
\end{center}\footnotetext[1]{\label{orgccd796c}Sedgwick, "Paranoid Reading and Reparative Reading." \emph{Touching
Feeling}.}\footnotetext[2]{\label{org6138d22}Butler, Judith. \emph{Bodies that Matter}.}\footnotetext[3]{\label{org69045c6}Muñoz, Jose Esteban. \emph{Disidentifications}.}\footnotetext[4]{\label{org2fa2be1}Spillers, Hortense. "Mama's Baby, Papa's Maybe".}\footnotetext[5]{\label{org0e22361}Prosser, Jay. \emph{Second Skins}.}

For example, chapter 2, "Text Analysis," explores the technology of
\emph{text analysis} and its mechanism of \emph{iteration}: a programmatic logic
of repetition in which text analysis programs execute the same action
across large collections of text, one word at a time, with the
attendant effect of supressing minor details in language forms in
favor of more frequent textual patterns. The chapter then takes Judith
Butler's concept of \emph{Gender Performativity} as a conceptual tool to
re-work iteration into a new text analysis practice for studying
gender in language.

\subsubsection{Communicating Across The Disciplines}
\label{sec:org63152f3}

In my research, disciplinary boundary-crossing aims to engage
knowledges that otherwise may not come into contact with one another.
Over the past two years, I have written three papers for primarily
Computer Science (CS) fields like Computational Linguistics and AI
Ethics. I have written two papers for the Humanities fields of Trans
studies and digital scholarly editing.

My CS papers include the aforementioned "What LLMs Can't See About
Gender," under review at \emph{FAccT 2027}, as well as two other papers:
"Some Myths About Bias: A Queer Studies Reading Of Gender Bias In
NLP," published in the \emph{Association for Computational Linguistics 2025
Proceedings}, and "Small Language Models and the Synthetic Middle,"
under review at \emph{NeurIPS 2026}. These papers draw from theorizing in
Queer studies and Trans studies in order to make recommendations to CS
fields. "Some Myths About Bias" draws from Queer Studies theorizing on
the binary to argue that bias is embedded in language in ways that
require more careful handling to avoid falling into the trap of binary
thinking. The other paper, "Small Language Models and the Synthetic
Middle," incorporate close reading into a novel, interdisciplinary
methodology to study overlap (the "synthetic middle") between
polarized discourses on gender and transgender rights in the USA.

My Humanities work includes the aforementioned paper in Trans studies,
"'My Body Feels Like it’s Coming Off': on Cis Embodiment in \emph{Love Is
Blind}," and another in the field of digital scholarly editing,
"Re-encoding Dominance: Queer Approaches To TEI Markup," published in
\emph{Digital Editing \& Publishing in the 21st Century}. This paper
examines how a language technology, TEI (Text Encoding Initiative), a
markup standard for encoding text-based cultural heritage materials,
might be read alongside theories from Queer studies and historical
work on the archive of slavery. 
\end{document}
